\documentclass[12pt]{article}
\usepackage[utf8]{inputenc}
\usepackage[T1]{fontenc}
\usepackage{hyperref}
\title{Yos Language Resources for K-12 Curriculum Development}
\author{Compiled by Claude Sonnet 5.5 (claude-sonnet-5-5) for Andr\'as Kornai}
\date{}
\begin{document}
\maketitle

\section{Introduction}
The code \texttt{yos} is not an active language. ISO 639-3 retired it, merging it into Zo (\texttt{zom}), under Change Request 2012-035, effective 2013-01-23. The stated reason was that ``Yos'' is the English plural of ``Yo'', a variant pronunciation of ``Zo'', and that its only other reference was an obscure mention in the 1983 Census of Myanmar. Glottolog lists it under ``Bookkeeping'', its category for retired codes. IANA marks the subtag deprecated, with \texttt{zom} as the preferred value.

\section{Primary Educational Resources}
None. There is no separate Yos language.

\section{Linguistic and Reference Materials}
\begin{itemize}
\item Glottolog entry for the retired code: \url{https://glottolog.org/resource/languoid/iso/yos}
\item IANA subtag record: \url{https://www.iana.org/assignments/lang-subtags-templates/yos.txt}
\end{itemize}

\section{Religious and Community Resources}
None.

\section{Digital and Audio Resources}
None.

\section{Curriculum Development Context}
Curriculum work should target Zo (\texttt{zom}) instead. Do not confuse \texttt{yos} with \texttt{yso}, the code for the Nisoish (Yi) languages. This dossier is intentionally minimal.
\end{document}